% Experimental Setup (Section 5) -- ML4H/PMLR Overleaf project.
% Requires: booktabs (tables), natbib (citations).
% Contains \label{sec:setup-eval}, referenced by sac.tex (cluster bootstrap).
% >>> NUMBERS IN TABLE 2 ARE FROM THE CURRENT GRADED FILES; regenerate from the
%     final run before submission. <<<
% Bib keys used: manes2024kqa, jeong2024olaph, abacha2017liveqa,
%   abacha2019medicationqa (last three: VERIFY in references.bib).

\section{Experimental Setup}
\label{sec:setup-exp}

\subsection{Datasets}
\label{sec:datasets}
We evaluate on three English medical question-answering benchmarks. \textbf{K-QA}
\citep{manes2024kqa} contains real patient questions whose gold answers are
decomposed by physicians into atomic statements, and ships with a physician-labeled
natural-language-inference set that we use to validate the factuality labeler.
\textbf{LiveQA} \citep{abacha2017liveqa} and \textbf{MedicationQA}
\citep{abacha2019medicationqa} are consumer-health and medication question sets; we
take them with the reference long-form answers collected in MedLFQA
\citep{jeong2024olaph}, the same medical long-form QA resource used by
\citet{cherian2024llm}. Each reference answer supplies the physician-verified
statements against which generated claims are graded. We pool the three datasets
within each model (Section~\ref{sec:setup-eval}); per-dataset breakdowns are in
Table~\ref{tab:data}.

\subsection{Models and confidence scores}
\label{sec:models}
The systems under test are two open instruction-tuned models,
\textbf{Llama-3-8B-Instruct} and \textbf{Mistral-7B-Instruct-v0.3}. For each
question the model generates a long-form answer, which we decompose into atomic
claims. The confidence score $s(c)$ is the model's own estimated probability that a
claim is true: we present the claim back to the generating model with a
\emph{True}/\emph{False} verification prompt and take the softmax over the two
answer-token logits,
\begin{equation}
s(c) \;=\; \frac{\exp(z_{\textsc{true}})}
{\exp(z_{\textsc{true}}) + \exp(z_{\textsc{false}})}.
\label{eq:ptrue}
\end{equation}
The score is thus intrinsic to the model under test and requires no external
resource at deployment. Before calibrating we require the score to be informative,
gating on $\mathrm{AUROC} \ge 0.7$ against the calibration labels; a score at chance
would still yield valid risk control but useless retention (Section~\ref{sec:sac}).

\subsection{Labeling: factuality and severity}
\label{sec:labeling}
Two labels are attached to each claim during calibration only. The \textbf{factuality
label} $y(c)$ is produced by a GPT-4o judge acting as a medical fact-checker,
grounded in the physician-verified reference statements of the corresponding
question: a claim entailed by the evidence is labeled true ($y=0$), one contradicted
by it a hallucination ($y=1$), and one neither entailed nor contradicted
unverifiable. Unverifiable claims are excluded from calibration and evaluation, so
all reported risks are computed over verifiable claims. We validate this labeler
against the K-QA physician NLI set and report agreement (accuracy and Cohen's
$\kappa$) in Table~\ref{tab:data}. % TODO: fill agreement number(s).

The \textbf{severity tier} $t(c)$ is assigned by a separate GPT-4o call under a fixed
rubric: a claim is \textsc{dangerous} if acting on it while false could cause direct
patient harm (dosing, contraindications, drug interactions, emergencies) and
\textsc{benign} otherwise (background facts, mechanisms, epidemiology, definitions).
% TODO: report severity-tier agreement against a physician-labeled subset (kappa),
% or soften the "physician-validated" claim in Related Work and Section 4.

\subsection{Calibration, evaluation, and confidence intervals}
\label{sec:setup-eval}
We evaluate with repeated calibration/test splitting. Within each pool we draw
$N_{\text{splits}} = 300$ random splits of verifiable claims; on each split we fit the
global threshold \eqref{eq:crc-threshold} and the per-tier thresholds
\eqref{eq:sac-threshold}, then measure, on the held-out claims, the dangerous-tier
risk, the true-claim retention, and (for reference) the marginal AUROC. The budgets
are $\alpha_{\textsc{dangerous}} = 0.05$, $\alpha_{\textsc{benign}} = 0.15$, and a
marginal baseline $\alpha = 0.10$. We report the mean over splits together with the
per-split \emph{violation rate}, the fraction of splits on which the realized
dangerous risk \eqref{eq:realized-risk} exceeds $\alpha_{\textsc{dangerous}}$.

Because claims from the same answer are correlated, the sampling unit for uncertainty
is the question, not the claim. We therefore form 95 percent confidence intervals by
a \textbf{question-level cluster bootstrap}: we resample whole questions with
replacement, recompute the fitted thresholds and the resulting tier risks on each
resample, and take the 2.5 and 97.5 percentiles across resamples. This is the
finite-sample counterpart to the claim-level exchangeability idealized in
Proposition~\ref{prop:sac}.

\subsection{Score ablation}
\label{sec:ablation}
To show that risk control is a property of the method rather than of any particular
score, we rerun the full calibration while holding the claims, labels, and tiers
fixed and swapping only $s(c)$ for: (i) a uniform random score, the noise floor;
(ii) the entailment probability of a natural-language-inference model; and (iii) the
length-normalized token log-probability of the claim. Because conformal risk control
is distribution-free, every score keeps the dangerous tier within budget; only
retention degrades, and it degrades most for the random score.

\begin{table*}[t]
\centering
\footnotesize
\setlength{\tabcolsep}{6pt}
\begin{tabular}{llrrrr}
\toprule
Model & Dataset & Questions & Claims & Verifiable & Dangerous halluc. \\
\midrule
\multirow{4}{*}{Llama-3-8B-Instruct}
  & K-QA          & 201 & 5{,}040  & 2{,}150 & 62 \\
  & LiveQA        & 100 & 2{,}131  & 953     & 48 \\
  & MedicationQA  & 425 & 10{,}201 & 4{,}072 & 286 \\
  & \textit{Pooled} & \textit{726} & \textit{17{,}372} & \textit{7{,}175} & \textit{396} \\
\midrule
\multirow{4}{*}{Mistral-7B-Instruct-v0.3}
  & K-QA          & 201 & 1{,}593 & 864     & 24 \\
  & LiveQA        & 100 & 901     & 455     & 16 \\
  & MedicationQA  & 666 & 5{,}357 & 2{,}669 & 106 \\
  & \textit{Pooled} & \textit{967} & \textit{7{,}851} & \textit{3{,}988} & \textit{146} \\
\bottomrule
\end{tabular}
\caption{Dataset statistics per model. \emph{Verifiable} counts claims labeled true
or hallucination (unverifiable claims are excluded from calibration and evaluation);
\emph{Dangerous halluc.} counts verifiable dangerous-tier hallucinations, the events
that determine statistical power. MedicationQA grading for Llama-3-8B was capped at
425 of 666 questions by a judge-API quota limit; because each model is calibrated and
evaluated independently, matched question counts across models are not required, and
the 396 pooled dangerous events for Llama already exceed the count needed for stable
estimates. Regenerate from the final run before submission.}
\label{tab:data}
\end{table*}
